%%%%%%%%%%%%%%%%%%%%%%%%%%%%%%%%%%%%%%
% LaTeX poster template
% Created by Nathaniel Johnston
% August 2009
% http://www.nathanieljohnston.com/index.php/2009/08/latex-poster-template/
%%%%%%%%%%%%%%%%%%%%%%%%%%%%%%%%%%%%%%

\documentclass[final]{beamer}
\usepackage[scale=1.24,orientation=portrait]{beamerposter}
\usepackage{graphicx}
\usepackage{ragged2e}
\usepackage[percent]{overpic}
\usepackage{lipsum}
\usepackage{bbm}

%-----------------------------------------------------------
% Define the column width and poster size
% To set effective sepwid, onecolwid and twocolwid values, first
% choose how many columns you want and how much separation you want
% between columns
% The separation I chose is 0.024 and I want 4 columns
% Then set onecolwid to be (1-(4+1)*0.024)/4 = 0.22
% Set twocolwid to be 2*onecolwid + sepwid = 0.464
%-----------------------------------------------------------
\newlength{\sepwid}
\newlength{\onecolwid}
\newlength{\twocolwid}
\newlength{\threecolwid}
\setlength{\paperwidth}{33.11in}
\setlength{\paperheight}{46.81in}
\setlength{\sepwid}{0.024\paperwidth}
\setlength{\onecolwid}{0.22\paperwidth}
\setlength{\twocolwid}{0.464\paperwidth}
\setlength{\threecolwid}{0.708\paperwidth}
%\setlength{\topmargin}{-0.5in}
\usetheme{confposter}
\usepackage{exscale}

%-----------------------------------------------------------
% Define colours (see beamerthemeconfposter.sty to change these colour definitions)
%-----------------------------------------------------------
\setbeamertemplate{caption}[numbered] % added this one myself.
\setbeamercolor{block title}{fg=LCSBgray,bg=white}
\setbeamercolor{block body}{fg=black,bg=white}
\setbeamercolor{block alerted title}{fg=white,bg=black}
\setbeamercolor{block alerted body}{fg=black,bg=white}

%-----------------------------------------------------------
% Name and authors of poster/paper/research
%-----------------------------------------------------------
\usepackage{tikz,tcolorbox}
\usetikzlibrary{shapes,calc}
\tcbuselibrary{skins}

\makeatletter
\newtcbox{\titlebox}[1][]{
	enhanced,
	overlay={
		\draw[#1,fill=#1](frame.south east)--+(0,.2)to[bend right]+(.2,-0)--cycle;},
	colback=#1,
	top=6pt,bottom=-2pt,left=2pt,right=2pt,
	boxrule=6pt,
	colframe=#1,
	sharp corners=south,
	colupper=white,
	fontupper=\bfseries\large
}
\newtcolorbox{myblock}[2][]{
	enhanced,
	left=2pt,
	right=2pt,
	colframe=#2,
	sharp corners=northwest,
	sharp corners=south,
	boxrule=6pt,
	colback=white,
	overlay={
		\def\myblock@tempa{#1}
		\ifx\myblock@tempa\@empty
		\else
		\draw[#2,fill=#2]($(frame.north west)+(.2pt,-.2pt)$)--+(.1,0)to[bend right]+(-0,-.1)--cycle;
		\node [
		anchor=south west,
		inner sep=0pt,
		outer sep=0pt
		]at(frame.north west){\titlebox[#2]{#1}};
		\fi
	},
	fontupper=\rmfamily
}
\makeatother

%-----------------------------------------------------------
% Useful user-defined commands
%-----------------------------------------------------------

% symbols for itemize enviroment
\newcommand{\bluebox}{\raisebox{1.5mm}{\tiny$\blacksquare$}\hspace{0.3cm}}
\newcommand{\bluestar}{\raisebox{1.5mm}{$\diamond$}\hspace{0.3cm}}
\newcommand{\bluebullet}{\raisebox{1.5mm}{$\bullet$}\hspace{0.3cm}}

% boxing text
\newcommand\mybox[2][]{\tikz[overlay]\node[fill=#1,draw=#1,inner sep=3pt, anchor=text, rectangle, rounded corners=2mm] {#2};\phantom{#2}}



% ----------------------------------------------------------------
% Main body of poster
% ----------------------------------------------------------------
\title{Your Title Here ...}
\author{Zuogong Yue and Jorge Gon\c{c}alves}
%\institute{Luxembourg Centre for Systems Biomedicine, University of Luxembourg}

% user's included packages here:
\usepackage{helvet}
\usepackage{times}
%\renewcommand{\familydefault}{\sfdefault}
\usepackage{algorithm}
\usepackage{algpseudocode}
\renewcommand\footnotesize{\small}


%-----------------------------------------------------------
% Start the poster itself
%-----------------------------------------------------------

% setup the top background/template
\usebackgroundtemplate{\includegraphics[width=\paperwidth]{poster-layout/top-blue.jpg}}

\begin{document}
\begin{frame}[t]

  % setup the bottom background/template
  \raisebox{-112cm}{
    \hspace*{-1.6cm}
    \begin{overpic}[width=1.05\textwidth]{poster-layout/bottom-blue.jpg}
      \put (4,14) {\color{white}{\large University of Luxembourg \hspace{12mm}$\bullet$  \hspace{12mm}Luxembourg Centre for Systems Biomedicine \hspace{12mm} $\bullet$ \hspace{12mm} www.uni.lu/lcsb \hspace{12mm} $\bullet$ \hspace{12mm} zuogong.yue@uni.lu }}
    \end{overpic}}

  \vspace{-110cm}  %<----- adjust this value to have right margin
                   % between body and title page
  \begin{columns}[t]

    % ----------------------------------------------------------------
    % FIRST COLUMN: the left half-page
    % ----------------------------------------------------------------
    % \begin{column}{0.3\sepwid}\end{column}    % padding to the left margin

    \begin{column}{\twocolwid}
      \setbeamercolor{block title}{fg=LCSBgray,bg=white}
      \vskip1ex
      \vspace{1cm}

      \begin{block}{\textsf{Introduction}}
        \large
        bla bla bla...
      \end{block}

      \vspace{1cm}
      \begin{block}{\textsf{Problem Formulation}}
        \large
        bla bla bla...

        This word is
        \mybox[LCSBred]{\color{white}{\textsf{highlighted}}} in red
        and in sans-serif fonts.
        This word is \mybox[LCSBblue]{\color{white}{highlighted}} in red.
      \end{block}

      \vspace{3cm}
      \begin{myblock}[Failure of discrete-time methods]{LCSBred}
        \large
        \vspace*{.5cm}
        \begin{center}
          \begin{minipage}[t]{0.95\textwidth}
            bla bla...
          \end{minipage}
        \end{center}

        \begin{figure}[htb]   %htbp!H
          \centering
          \includegraphics[width=.8\textwidth]{poster-layout/top-red}
        \end{figure}
        \vspace*{.5cm}
      \end{myblock}

    \end{column}% END of the left half-page


    % ----------------------------------------------------------------
    % SECOND COLUMN: the right half-page
    % ----------------------------------------------------------------
    % \begin{column}{0.5\sepwid}\end{column}	    % padding to the left margin

    \begin{column}{\twocolwid}
      \setbeamercolor{block title}{fg=LCSBgray,bg=white}
      \vspace{1.55cm}

      \begin{block}{\textsf{Network inference}}
        \large
        bla bla...

        \begin{itemize}
        \item[\bluebullet] item 1
        \item[\bluebullet] item 2
        \item[\bluebullet] item 3
        \end{itemize}

        \vspace*{2.7cm}
        \begin{myblock}[\textsf{Algorithm}]{LCSBblue}
          \vspace*{.3cm}
          \begin{minipage}[t]{.95\textwidth}
            \begin{algorithmic}[1]
              \State bla bla...
              \State bla bla...
            \end{algorithmic}
          \end{minipage}
          \vspace*{.3cm}
        \end{myblock}

      \end{block}


      \begin{block}{\textsf{Conclusions}}
        \large
        bla bla ...
      \end{block}

      \vspace{1cm}
      \begin{block}{\textsf{References}}
        { \small
          \begin{thebibliography}{1}

          \bibitem{GW}
            \justifying{
              J.~Goncalves and S.~Warnick.
              ``Necessary and Sufficient Conditions for Dynamical Structure
              Reconstruction of LTI Networks''.
              \emph{Automatic Control, IEEE Transactions on},
              \textbf{53}(7):1670--1674, 2008.
            }

          \bibitem{GN}
            \justifying{
              S.~Gibson and B.~Ninness.
              ``Robust maximum-likelihood estimation of multivariable dynamic
              systems''. \emph{Automatica}, \textbf{41}(10):1667-1682, 2005.
            }

            % \bibitem{Wipf}
            %   \justifying{
            %   D.P.~Wipf, B.D.~Rao.
            %   ``Sparse Bayesian learning for basis selection''.
            %   \emph{Signal Processing, IEEE Transactions on},
            %   \textbf{52}(8):2153-2164, 2004.
            % }

          \end{thebibliography}
        }
      \end{block}

    \end{column}% END of the right half-page
    % \begin{column}{\sepwid}\end{column}      % padding to the right margin

  \end{columns}
\end{frame}
\end{document}
